\documentclass[11pt,a4paper]{article}
\usepackage[T1]{fontenc}
\usepackage[utf8]{inputenc}
\usepackage{lmodern}
\usepackage[margin=25mm]{geometry}
\usepackage{amsmath,amssymb,booktabs,tabularx,xurl,xcolor,hyperref}
\hypersetup{colorlinks=true,linkcolor=blue!50!black,urlcolor=blue!50!black,
 pdftitle={Relation Ownership and Bag Preserving Transfer},pdfauthor={Mettapedia Working Notes}}
\setlength{\emergencystretch}{3em}
\setlength{\parindent}{0pt}
\setlength{\parskip}{5pt}
\newcommand{\code}[1]{\nolinkurl{#1}}
\title{Relation Ownership and Bag Preserving Transfer}
\author{Mettapedia Working Notes}
\date{1 November 2026}
\begin{document}
\maketitle
\begin{abstract}
Eighteen Lean modules establish independent semantic contracts for relation-column
ownership, retain/release receipts, dead-ID reuse and row take-over. The
bookkeeping algorithm updates each column occurrence while maintaining one
external hold per used ID. Transfer preserves the full occurrence bag, even
when duplicate copies remain as rows or the target gives its representation
back during the pass. Publication composes this bag law with the existing
captured-reader protocol. Separate root owners expose descriptor and snapshot
retention, and local bridges connect semantic movement and zero-crossing
counts to existing cost models. Stable renumbering preserves decoded rows,
copied use counts and side-index candidates; original logical coordinates
connect physical removal marks to the bag law. A UInt32 realization preserves
counts and receipts under a whole-row capacity condition. Unique binary
relations inherit a stronger bound from their finite numbering; general tuple
tables do not. Native pointer updates, enforcement of the range conditions,
root discovery and publication isolation remain implementation obligations.
An incremental publication worker preserves its compensated bag at every
private step. Snapshot root enumeration and a separate native-metadata lease
model compose term collection with descriptor retirement. Pinned-executable
allocation-failure controls distinguish recovery, fatal termination and a
deliberately broken publication step. An additional matching-admission law
exposes a false absence certificate for open compact facts in the pinned
runtime; the expanded native gate is not green.
\end{abstract}

\section{Column multiplicity and external ownership}

Let $R$ be the list of live relation entries, each entry a list of stored IDs.
The reference demand for ID $i$ is
\[
 D_R(i)=\sum_{r\in R}\operatorname{count}_i(r).
\]
This is the existing \code{ArenaHeldResources.logHolds} interpretation, reused
without redefining it to be the target bookkeeping algorithm.

\code{RelationColumnOwnership} stores a counter function $u$ and a finite
external-ownership set $H$. Its invariant is
\[
 u(i)=D_R(i),\qquad i\in H\ \Longleftrightarrow\ u(i)>0.
\]
Insertion walks the columns, increments each counter and takes a hold only
when the previous counter was zero. Removal walks them in the same manner,
decrements and releases only when the previous counter was one. The proofs
establish the invariant for insertion and membership-checked removal.

A represented row supplies enough demand for every decrement, including
repeated occurrences of the same ID. The safe-removal theorem proves this
at each intermediate step, not just for the final subtraction. For example,
the row $[7,7]$ contributes two internal uses and one external hold. Removing
it releases the hold only after the second decrement. If another row
$[7,9]$ remains, ID 7 stays held.

The unique-insert wrapper refuses a row already represented in the dense
relation. This is separate from the surrounding space's bag semantics:
additional copies can remain in plain rows. A presence bit is therefore
neither a column-use count nor a complete occurrence count.

\section{Callbacks are occurrences too}

An external reference counter observes callback invocations. Equality of
owned sets alone would not detect two redundant retain callbacks.
\code{RelationHoldReceipts} executes the zero-crossing receipt algorithm and
proves that its acquired and released ID lists have no duplicates. Membership
is characterized exactly:
\begin{align*}
 i\in\mathsf{acquired}(u,r)
   &\Longleftrightarrow i\in r\ \land\ u(i)=0,\\
 i\in\mathsf{released}(u,r)
   &\Longleftrightarrow i\in r\ \land\ u(i)=\operatorname{count}_i(r),
\end{align*}
where the release theorem requires sufficient initial counts. These lists
are sublists of the input columns. They consequently emit exactly one
callback for each required ownership transition and at most one per column.

The ownership-change theorems identify acquisition with newly owned IDs and
release with IDs no longer owned afterward. The controls reject a naive
per-column callback sequence $[7,7]$ for a newly inserted repeated-column row:
the correct retain receipt is $[7]$.

The model uses natural counters. Both inspected C stores use
\verb|uint32_t holds[]| and unchecked increment/decrement operations. The
logical safe-removal result supplies a no-underflow premise when the native
representation invariant holds. The following finite-word results make the
overflow condition precise; they do not demonstrate an allocated overflowing
runtime workload.

\section{Finite words and relation-specific capacity}

\code{RelationWordOwnership} implements the updates with Lean's actual
\code{UInt32} arithmetic: hold tests the old word and increments it; unhold
decrements first and tests the new word. Interpretation takes each word's
natural value and retains the ownership set. For an inserted row $r$, the
necessary range premise used by the refinement is
\[
 \forall i,\quad u(i)+\operatorname{count}_i(r)<2^{32}.
\]
Under that premise, the complete word update equals the natural-count
update, including the ordered acquire receipt. Removal of a represented row
inherits the earlier no-underflow proof and preserves the release receipt.
The resulting represented store remains valid.

The executable preflight counts all occurrences of an ID in the row before
performing any update. A second version compares the representable addition
against \verb|UINT32_MAX - old|; it is proved equivalent to the natural
condition and does not evaluate an overflowing sum. Preparation refusal
publishes no updated counter state. The inspected C code does not currently
perform this preflight; the result is a permission law for such admission.

Checking room for one use against the old counter is insufficient for a
row containing an ID twice. A control starts ID 7 at $2^{32}-2$ and inserts
$[7,7,9]$: unchecked increments wrap 7's count to zero while its ownership
set still contains 7. The checked preflight refuses that row but accepts
$[7,9]$. Decrementing zero likewise wraps, rather than acting as truncated
natural subtraction.

\code{RelationOccurrenceBounds} bounds every ID's count by total column
volume and then by arity times fact count. This is sufficient but not always
sharp: different IDs can each fit while their combined volume exceeds one
counter's range.

For unique binary pairs, \code{RelationBinaryUseBound} proves a stronger law.
With $n$ numbered IDs, a fixed ID has at most $n$ distinct partners in each
column, so its use count is at most $2n$, independent of the number of pairs.
The numbering constraints
\[
 4n\leq3s,\qquad s\leq2^{31}
\]
imply $n<2^{31}$. A fresh pair therefore satisfies the whole-row word
admission condition. This licenses omission of a separate counter preflight
\emph{if} the native representation maintains pair uniqueness, correct counts
and those numbering bounds. A repeated-pair bag fails the bound, which is
why extra bag copies staying outside the unique relation matters.

Higher arities do not inherit that argument. The checked symbolic family
\[
 [x,x,x,x,x,x,x,y_j],\qquad 0\leq j<613\,566\,757,
\]
has distinct eight-column rows and at most $613\,566\,758$ IDs. Both its row
count and ID count fit a three-quarter-loaded $2^{30}$-slot table, but $x$
has $4\,294\,967\,299>2^{32}$ column uses. The proof instantiates general
list/count laws; it does not allocate or enumerate that runtime workload.
It shows that uniqueness and these hash-load bounds alone cannot discharge
the tuple counter's range obligation. A further control compares release
receipts after that count is encoded in a word: removing seven uses of $x$
and the sole use of another ID releases both in the wrapped worker, whereas
the natural-count worker releases only the other ID. Overflow can therefore
change ownership observations, not merely a private numerical statistic.

\section{Dead dictionary entries and identifier reuse}

An unused number may remain in a dictionary without retaining the term it
once named. \code{RelationIdReuse.reserve} reuses an existing raw-ID entry or
appends a new one. It proves that the returned index names the requested ID,
old indices are preserved, and distinct dictionaries stay distinct.
The state, dictionary and lookup share one admitted identifier namespace;
cross-universe reuse or reset requires its own revalidation contract.

The dictionary contains IDs, not cached payload meanings. Reading through a
reserved number consults the current lookup. Thus a dead entry naming raw
ID 7 can correctly name a newly allocated payload after that ID is reused.
A retained old-payload cache would be wrong. Missing lookups remain explicit
\code{none} values in the denotation; they are not filtered into apparent
absence of a fact.

The column invariant proves that a zero-use ID occurs in no live relation
row. Changing its meaning therefore preserves that relation's existing rows.
This is not yet global permission to reuse it: another relation, a caller,
a descriptor or a captured reader may still protect it. The checked reuse
operation requires absence from a supplied complete protected set and
preserves every lookup in that set.

Positive controls remove the old repeated-column row, retain its dead number,
change the raw ID's meaning and reinsert a row that observes the new payload.
Negative controls show both stale-payload reuse and a stale side entry
inventing a fact. Correct native side indexes must describe the same live
entries as the counters; the dictionary theorem alone does not establish
that array/index correspondence. The next section supplies the finite-list
renumbering and candidate-index part of the obligation.

\section{Stable renumbering and side-index candidates}

For an old dictionary $D$ and live mask $b$, define
\[
 \rho_b(i)=\operatorname{count}_{\mathsf{true}}(\operatorname{take}_i b).
\]
\code{StableNumberCompaction.pack} scans the mask and dictionary together,
retaining live entries in their original order. With equal mask and dictionary
lengths, a live coordinate satisfies
\[
 \mathsf{pack}(b,D)[\rho_b(i)]=D[i].
\]
Lookups are optional, so invalid coordinates are not silently replaced by a
default value. The map is strictly increasing on live coordinates, satisfies
$\rho_b(i)\leq i$, and covers every compact coordinate without gaps. Sorted
side arrays therefore remain sorted. For bit rows the destination word is
never later than the source word. That last bound is a necessary condition
for the inspected in-place loop, not a complete proof of its mutable writes:
the current source word must still be captured before it is cleared.

\code{RelationNumberingCorrespondence} independently decodes numbered rows
through $D$. It builds $b$ from positive column-occurrence counts, packs $D$
and maps every row coordinate through $\rho_b$. For a valid dictionary and
in-range columns, the entire decoded row sequence is unchanged, the output
is valid and every retained number is used. The compact dictionary contains
exactly the IDs named by live columns.

The copied hold-counter array also agrees with the relocated row counts.
For live $i$,
\[
 D_{\mathsf{rows}_{\mathrm{after}}}(\rho_b(i))
   =D_{\mathsf{rows}_{\mathrm{before}}}(i).
\]
This uses injectivity on live coordinates; it is not valid for arbitrary
dead coordinates. A checked control shows a dead coordinate's rank aliasing
a live one, which is why dead side entries cannot be relocated as facts.

An exact side-query law preserves the original row positions selected by a
column and key, after relocating the key and columns together. Repeated
equal rows retain separate positions. This specifies the required content
of rebuilt hash tables and side chains without claiming their C mutation
code has been verified. A control with dictionary $[8,7,6,9]$ and rows
$[[1,1],[3,1],[1,1]]$ compacts to dictionary $[7,9]$, rows
$[[0,0],[1,0],[0,0]]$ and hold counts $[5,1]$.

\section{The roots have different owners}

The inspected runtime separately retains a kept relation's head and the
constructor IDs in its shared skeleton. Removing its last fact column does
not automatically retire that descriptor. Accordingly,
\code{RelationRootOwners} separates
\[
 H_{\mathrm{columns}},\quad H_{\mathrm{descriptor}},\quad
 H_{\mathrm{captured}}.
\]
Their union is the direct root support of the declared owners. It is embedded
in the existing \code{ResourceOwnership} finite-heap model, whose reachable
footprint includes shared and cyclic descendants. Reuse is constrained by
that complete footprint rather than merely by the direct roots.

Under root-allocation validity, every live column resolves to a payload.
Protected reuse preserves live relation denotation, descriptor reads and
captured-owner reads. A separate bridge derives the captured roots from the
existing reader's saved row view and proves preservation of the values read
from its remaining suffix.

These root sets describe reachability support, not callback multiplicity.
The exact callback accounting is the separate receipt result above.
Nor does this ownership split prove that retaining an empty descriptor is
the best cache policy. A stronger policy that no metadata IDs remain after
the last fact requires retirement or reconstruction of that descriptor.
The current column theorem must not be advertised as proving that stronger
property.

\subsection{A snapshot needs both term roots and live metadata}

\code{RelationCapturedRoots} models the inspected callback: enumerate the
head, the constructor IDs of non-hole skeleton nodes, all declared columns,
and every materialized record-cache ID. Its independent reader observes the
complete supplied fields. Under the layout equations
\[
 |\mathsf{columns}|=\mathsf{count}\cdot\mathsf{arity},\qquad
 |\mathsf{records}|=\mathsf{count},
\]
the callback list contains exactly the IDs required by that reader.
Collection retaining these roots preserves every field and every reachable
path. Filling an empty record-cache slot preserves the layout and adds the
record's ID to the next enumeration. A root list captured before that fill
is not sufficient for a later collection.

The callback itself also reads native memory.
\code{RelationSnapshotMetadata} separates the native-address heap from the
term-ID heap. Its loader requires the snapshot allocation, its owning
universe and its separately allocated skeleton. Keeping every term alive
does not reconstruct a freed skeleton or universe object. A prepared census
checks the universe address of every provider and refuses the entire census
if any provider fails to load; it never turns a failed provider into an
empty root list. Completeness here is relative to the supplied provider
registry. Discovering every native owner remains a separate obligation.

The model reuses \code{ResourceOwnership.AcquiredUse}: detaching a
descriptor's current slot leaves an already admitted reader's lease valid.
The final reader can extract the retired descriptor exactly once. The
composition theorem preserves the full snapshot readout after descriptor
detachment, native-metadata collection and term collection. It retains
outside roots in both heaps. The theorem requires valid layouts, allocated
term roots, the same admitted universe and live metadata dependencies;
native pointer access, provider synchronization and universe-reset protocols
must realize those conditions.

Controls distinguish a missing cache entry from a dangling materialized
record, reject a foreign-universe provider, and show why keeping only column
IDs is insufficient. Another control frees the skeleton while preserving all
term lookups: the metadata loader correctly fails.

\subsection{The provider registry has a semantic interpretation}

\code{RelationProviderRegistry} reuses \code{RegistryIndexCost}'s existing
slot array, coherent index, append and swap-removal operations. For a current
callback enumeration $E$, its root census is
\[
 C(R,E)=\bigcup_{p<R.\mathsf{len}} E(R.\mathsf{slot}(p)).
\]
Index coherence proves that this is exactly the union over registered keys.
Adding key $k$ changes the census to $E(k)\cup C(R,E)$; adding it twice does
not add another registry slot. Removing $k$ is equivalent to enumerating the
old registry with only $k$'s enumeration replaced by the empty set. It is
\emph{not} generally $C(R,E)\setminus E(k)$, because other providers may
still name those IDs. The swap-removal proof retains the moved last slot's
roots at its new coordinate.

Composing that census law with the captured-root theorem preserves a
registered snapshot's readout across collection. The enumeration is evaluated
at collection time; registration does not cache its old result. A control
adds a materialized root to a registered provider and shows that a stale
census would miss it. The key carrier is natural numbers naming complete
callback/context identities within one admitted universe. Encoding those
identities and enrolling every external owner remain native obligations.

\section{Transfer preserves the full bag}

\code{RelationBagTransfer} models one target relation with unique dense entries
and rows into which it may give its representation back. The source is an
occurrence list. Each offer has one of four dispositions: refuse, store,
give back and refuse, or store and then give back. An unavailable relation
does not accept subsequent stores.

The pass records one Boolean decision per original source occurrence. A
successful offer removes that occurrence from the source; a refusal retains
it. Target give-back does not alter the original decision coordinates. For
every action schedule,
\[
 [\mathsf{remaining}] + [\mathsf{dense}_{\mathrm{after}}]
   + [\mathsf{spilled}_{\mathrm{after}}]
 = [\mathsf{source}] + [\mathsf{dense}_{\mathrm{before}}]
   + [\mathsf{spilled}_{\mathrm{before}}],
\]
where brackets denote multisets. This is a full bag equality, not equality
of supports or total sizes. The dense list remains duplicate-free whenever
it started so. Framing an unaffected bag preserves the equality, providing
the local composition law for other relations.

For source rows $[7,7,9]$ and an empty target, storing every admissible row
leaves $[7]$ as plain rows and dense entries $[9,7]$. Deleting the refused
second copy would lose an occurrence. If the first accepted entry is then
given back, it remains as a spilled row while later source occurrences stay
plain. Both cases are checked controls.

\code{RelationTransferCoordinates} converts the successful mask into strictly
increasing, bounded original logical positions. Independent deletion by
those positions equals the mask specification. If give-back appends rows,
the original positions leave that suffix untouched: their bounds refer to
the original prefix, not a newly shortened list.

A second worker marks a physical live-mask, advancing its logical coordinate
only at a live row. Decoding its result equals logical deletion; settling the
marks leaves that readout unchanged. Composing these laws with transfer proves
full bag preservation for the physical-mask representation and the give-back
suffix. For physical rows $[99,7,98,9,7]$ with live mask
$[0,1,0,1,1]$, removing logical position one removes 9, not the physical row
at position one. Controls distinguish that error and keep repeated 7s as
separate occurrences.

Native mark selection and array writes must realize these workers. The
one-relation framing law does not itself verify every mutable multi-head
pass or every native index update.

\section{Publication and captured readers}

Planning is private. Recoverable preparation refusal leaves the public
store unchanged, while successful publication installs the complete plan.
The copy-before-delete interval is not a public semantic state: publishing
one newly copied 7 while retaining both old copies would temporarily expose
three occurrences instead of two. A checked negative control exhibits that
difference.

\code{RelationTransferPublication} makes that interval executable in the
model. A draft holds processed original rows, their acceptance mask, pending
original rows and the target. Each step offers one pending row and appends
its actual disposition. The compensated logical store includes only refused
processed rows, all pending rows and the target. Its bag is invariant after
each step, even though the physical store still contains accepted source
copies. Any finite page preserves the invariant, and
\[
 \mathsf{advance}(a+b,d)=
 \mathsf{advance}(b,\mathsf{advance}(a,d)).
\]
One row is consumed from the worklist per step, so a budget at least the
original row count completes the model worker. This is a bound on abstract
offers, not on native instruction cost.

A session distinguishes available, working and halted states. A working
session returns \emph{busy}, not an empty bag; halted work returns
\emph{terminated}, not the old value. Publication accepts only completed
work. Every value returned from a consistent session is the original bag.
These distinctions specify the required isolation contract; they do not
assert that the native take-over loop is currently resumable.

Publication uses the existing detach-before-edit reader protocol.
\code{published_bag_and_reader} proves both that the newly published bag
equals the original bag and that the old reader's remaining sequence is
unchanged. With the example above, the new physical enumeration is
$[7,9,7]$. A reader already at position one must retain its old suffix
$[7,9]$; adopting the new revision without capture would instead read
$[9,7]$.
The incremental worker's \code{published_bag_and_captured_sequence} theorem
establishes the same pair of observations after any admitted completed session.

Native realization must protect the copy/delete window and enumerate the
captured roots. The inspected \code{cetta_malloc} and \code{cetta_realloc}
abort on allocation failure. A fatal abort after partial physical work is
not a recoverable refusal or a rollback certificate. Conversely, raw
allocation failures during preparation or relation insertion can follow
refusal/give-back routes. The obligation map keeps those dispositions distinct.

\subsection{Pinned native failure controls}

The r82 executable named in the qualification note was exercised with three
source shapes: nested facts in a tuple table, diagonal binary facts in array
rows, and binary star facts in an actual bit row. Reviewed call-site probes
simulate one NULL allocator result before the allocator runs, retaining the
old allocation in the realloc cases. Registers are deliberately changed in
these private child processes; this differs from the observation-only route
coverage probes. Executable, source and harness identities are checked.

Across the three shapes, thirteen recoverable-failure controls preserve the
complete reference observations. They cover initial and growing move storage,
the taken-position array, tuple/hold growth and binary hold growth. The latter
growth failures force give-back of 64 entries for the tuple/diagonal controls
and 63 for the star. The row-kind receipt distinguishes the actual array and
bit cases. Three snapshot-detachment allocation failures terminate with
SIGABRT. Three deliberate mutants skip source-row removal and are rejected
by an answer-bag mismatch. These are nineteen reached controls, not exhaustive
allocation-failure or concurrent-interleaving coverage. Fatal termination
does not establish rollback.

Retained-reader controls additionally keep one or seventeen work handles
across source requests, deliver a prefix, remove and replace the source
facts, and then finish every old reader. Sixteen combinations of nested or
flat facts, reader count, plain or eager relation storage, and disabled or
eager reclamation preserve the independent original-bag, post-state and
outside-owner oracles. The layout probe follows the frozen executable, not
the newer checkout's added lifetime-hold field.

Bounded tracing selects the first collection for each universe, active
snapshot-context set and record-cache generation. Across the four enabled
collection probes, 70 selected live collections invoke 582 checked snapshot
callbacks. Every selected census visits all its registered snapshot providers,
and each callback's exact root sequence agrees with its live fields. This
does not claim exhaustive interleaving coverage. The disabled-collection
controls provide no positive collection witness. A further fault at provider
registry growth reaches the underlying allocator and terminates with SIGABRT;
the defensive false-return branch is not credited as exercised.

\section{Existence certificates must include compact open facts}

The final observer audit found a separate failure in the pinned r82 engine.
With only \verb|(seen0 $open)| in an eagerly compacted set, ordinary
\verb|(match S (seen0 b) t)| finds an answer. Yet the native empty-collapse
test reports absence. A case-based insert-on-absence guard consequently adds
an extra ground fact. The direct match and its optimized existence observer
disagree; this is not explained by answer ordering.

The reviewed native trace returns false from
\code{space_contains_exact}, true from
\code{space_contains_only_exact_atoms}, and then false with applicability
set from \code{space_match_exists_ground_exact}. The inspected exactness
summary accounts for the ordinary row array, which is insufficient when
compact storage contains a variable-bearing fact. The reproducer survives
the source-variant checks and is reduced to one stored open fact and one
ground query.

\code{MatchExistenceAdmission} reuses the mixed store of
\code{ShapeHeadRouting}. It proves two different permissions. A positive
identity hit establishes existence when the query matches itself. A negative
hit needs a complete certificate that every matching logical candidate would
be that identical fact. An executable whole-store classifier folds both
ordinary and compact components and satisfies this contract when its local
classifier is sound. Equivalent representations produce the same summary.
The query-specific certificate may be more permissive than global groundness;
facts that provably cannot match need not prevent the optimization.

A one-column control interprets \code{none} as a fresh variable and
\verb|some n| as a ground atom. An empty ordinary row array incorrectly
certifies absence when the compact component contains \code{none}. The
complete summary declines that miss; a present ground fact still allows a
positive hit even alongside the variable. A final control shows the wrong
absence changing the post-effect bag size from one to two. These are checked
permission laws and counterexamples, not an applied C repair.

A later provenance-backed frozen successor candidate keeps the open first fact in ordinary
rows when its head subsequently takes a relation. The expanded existence
cases then agree, as do the complete joined qualification packages. This is
an observed repair in the candidate artifact; the work here neither changed
the C source nor treats the reconstructible build receipt as a cryptographic
ELF-to-source proof.

\subsection{The coordinate helper has a smaller observation domain}

The HE empty-case path uses a separate coordinate-form helper behind
\code{space_match_exists_flat_linear_view64}. Its answer is not automatically
a whole-space answer. \code{MatchExistencePartition} proves the required
composition:
\[
 \mathsf{exists}_{R+C}(q)=
 \mathsf{exists}_{C}(q)\lor\mathsf{exists}_{R}(q),
\]
where $C$ is compact storage and $R$ ordinary rows. A compact hit finishes
the Boolean observation. A certified compact miss licenses a row-only answer.
If compact matching is undecided, materializing the relevant head preserves
the whole bag and leaves no matching compact fact, after which rows are the
complete domain. Unrelated compact heads may remain stored.

Optional evidence keeps unknown distinct from false. The merge accepts any
positive witness, returns false only from two negative domain certificates,
and otherwise remains unknown. The model proves that an admitted merged
answer is correct. It also records the asymmetric mutation laws for a fixed
query: a positive witness survives additions, while absence survives
restrictions. Arbitrary insertion can invalidate absence and arbitrary removal
can invalidate presence.

Pinned r82 and successor-candidate controls exercise twelve HE scenarios:
ground, open and unrelated-open stores; plain and eager storage; direct and
explicitly materialized observation. Direct traces distinguish compact hit,
compact miss followed by rows, and settle followed by rows. Materialized
controls take no specialized view-existence call and supply the same fixed
answers. Eight negative harness controls reject cross-call, missing, late,
refused and incomplete evidence. These sequential caller-chain receipts do
not establish concurrent stability or source-to-binary correspondence.

\section{The semantic quantities entering the cost models}

\code{RelationOwnershipCostBridge} partitions the actual acquire receipt into
IDs absent from the dictionary and dead IDs already in it. Their total is
bounded by row arity; revived IDs are bounded by the dictionary's dead-ID
set. These facts make the insertion clamps in \code{RelationNumberingCost}
identities when its dead-count field denotes that set. Release receipts are
bounded by both row arity and prior live ownership.

For transfer, successful moves are bounded by independently counted eligible
store candidates. If the pending estimate covers those candidates, the
movement clamp in \code{RelationTransferCost.pass} is also an identity. The
remaining-row count agrees with subtracting the actual successful mask count.
The current imported cost model includes two row scans and the moves, allowing
five units per counted promotion row in its abstract work bound.

These are local interfaces between independently proved semantics and cost
accounting. Full native scheduling, room growth, failed reservations, mark
probes and machine time need their own trace and representation maps.

\section{Implementation obligations}

\begin{center}
\small
\begin{tabularx}{\linewidth}{@{}>{\raggedright\arraybackslash}p{.27\linewidth}>{\raggedright\arraybackslash}X>{\raggedright\arraybackslash}X@{}}
\toprule
Runtime seam & Checked semantic contract & Remaining native obligation\\
\midrule
\code{dense_hold}, \code{number_hold} & Word/count refinement, receipts and binary use bound & Preserve premises; enforce tuple bounds; execute callbacks\\
\code{dense_unhold}, \code{number_unhold} & Safe decrements and exact release & Live-entry/index correspondence\\
\code{dense_number_or_add}, \code{number_of} & ID reservation and current lookup & Array realization and failed reservation metadata\\
\code{dense_renumber}, \code{tuples_renumber} & Stable lookup, counts, sortedness and side candidates & Hash/side rebuilds and in-place bit writes\\
\code{space_dense_list} and snapshot root callbacks & Root enumeration, metadata leases and scoped census & Complete registration, native leases and synchronized updates\\
Provider add/remove & Coherent indexed census, shared roots and swap-removal & Callback/context encoding and complete enrollment\\
Exact existence & Whole-domain and partitioned negative-certificate laws & Include compact non-exact facts or retain the caller's compact-domain certificate\\
\code{space_dense_take_over_rows} & Stepwise bag conservation and original-coordinate deletion & Eligibility, physical realization and isolation\\
\code{space_dense_settle_head} & Give-back preserves the target bag & Exact skeleton reconstruction and ownership handoff\\
Pinned readers and removal marks & Bag, logical/physical mark bridge and retained readout & Pointer lifetime, detachment and array mutation\\
\bottomrule
\end{tabularx}
\end{center}

The eighteen mathematical modules and \code{RelationCorrespondenceAudit} pass the
targeted build and strict warning-as-error checks. Declaration manifests cover
1,216 checked bodies, with nineteen nonlogical recursion companions listed
separately. There are no admissions or added axioms; primitive dependencies
are confined to \code{propext}, \code{Classical.choice} and \code{Quot.sound}.
This is a model-level correspondence package and obligation map, not a
whole-C verification or a claim of information-minimal storage.
\end{document}
